% Generated draft: requires mathematical and source review.
\documentclass[11pt]{article}
\usepackage[margin=25mm]{geometry}
\usepackage{fontspec}
\setmainfont{Latin Modern Roman}
\usepackage{amsmath,amssymb,mathtools,unicode-math}
\setmathfont{Latin Modern Math}
\usepackage{enumitem,needspace,xurl,hyperref}
\setlength{\parindent}{0pt}
\setlength{\parskip}{5pt}
\setlength{\emergencystretch}{3em}
\newcommand{\N}{\mathbb{N}}
\newcommand{\Z}{\mathbb{Z}}
\newcommand{\Q}{\mathbb{Q}}
\newcommand{\R}{\mathbb{R}}
\newcommand{\C}{\mathbb{C}}
\newcommand{\F}{\mathbb{F}}
\newcommand{\A}{\mathbb{A}}
\newcommand{\PP}{\mathbb{P}}
\newcommand{\E}{\mathbb{E}}
\newcommand{\eps}{\varepsilon}
\newcommand{\dd}{\,\mathrm{d}}
\begin{document}
\section*{UnsolvedMath ID: 4200009}
\textbf{Source problem code:} catalogue entry
\subsection{Definitions and mathematical statement}
A Hamiltonian flow on a symplectic manifold $(M,\omega)$ is generated by a Hamiltonian $H$ through $\iota_{X_H}\omega=dH$; the sign convention has no effect on the mixing question. A smooth invariant $d$-torus is a smooth embedding $i:\mathbb T^d\to M$, where $\mathbb T^d=\mathbb R^d/\mathbb Z^d$, whose image is invariant under the flow $\Phi^t$. Its induced flow is $\psi^t=i^{-1}\circ\Phi^t\circ i$.

For an invariant Borel probability measure $\mu$ on $\mathbb T^d$, strong mixing means
\[
\lim_{t\to+\infty}\mu((\psi^t)^{-1}(A)\cap B)=\mu(A)\mu(B)
\]
for every pair of Borel sets $A,B$. Does there exist a Hamiltonian system with a smooth invariant torus for which the induced dynamics is mixing, especially with $d\geq3$? In the usual smooth-volume version one requires $\mu$ to have a smooth strictly positive density on the torus. The source does not explicitly state this measure hypothesis, so it must be fixed when interpreting the question; allowing a point mass at a fixed point would give a vacuous mixing example.

The supplied record states that no such mixing occurs on a two-dimensional torus, in the intended nontrivial class, and its literature summary reports analytic Hamiltonian examples with mixing restricted dynamics on invariant tori in dimension at least three. These are source claims; no construction, citation for the resolution, or additional hypotheses on the invariant measure are supplied. No Lagrangian condition on the invariant torus is imposed by the statement.
\subsection{Short English statement}
Can the dynamics of a Hamiltonian system be strongly mixing when restricted to a smooth invariant torus? The source reports higher-dimensional analytic examples and a two-torus obstruction, but does not specify the invariant measure or cite the resolution.
\subsection{Sources}
\begin{itemize}
\item Oliver Knill, \emph{Open problems in Hamiltonian dynamics} (2000), item 9, ``Kolmogorov problem,'' \url{https://people.math.harvard.edu/~knill/seminars/intr/}; accessed 2026-07-29, HTML extraction. The record attributes the question to Andrey Kolmogorov.
\item Local registry \texttt{data/unsolvedmath/problems.json}, permanent UnsolvedMath ID $4200009$, source code catalogue entry, \emph{Kolmogorov mixing-torus problem}. Dataset attribution: Ulam AI and the UnsolvedMath Contributors.
\item The supplied SOLVED-IN-LITERATURE summary gives no author, title, date, or publication identifier for the asserted analytic examples or the two-dimensional obstruction.
\end{itemize}
\end{document}
